\documentclass[10pt,a4paper]{amsart}
\usepackage{amsmath,amssymb,amsthm,graphicx}
\usepackage{verbatim}
\newenvironment{reference}{\comment}{\endcomment}
%\newenvironment{reference}{}{}
\newenvironment{slogan}{\comment}{\endcomment}
\newenvironment{history}{\comment}{\endcomment}

% For commutative diagrams we use Xy-pic
\usepackage[all]{xy}

% We use 2cell for 2-commutative diagrams.
\xyoption{2cell}
% Theorem environments.
%
\theoremstyle{plain}
\newtheorem{theorem}[subsection]{Theorem}
\newtheorem{proposition}[subsection]{Proposition}
\newtheorem{lemma}[subsection]{Lemma}

\theoremstyle{definition}
\newtheorem{definition}[subsection]{Definition}
\newtheorem{example}[subsection]{Example}
\newtheorem{exercise}[subsection]{Exercise}
\newtheorem{situation}[subsection]{Situation}

\theoremstyle{remark}
\newtheorem{remark}[subsection]{Remark}
\newtheorem{remarks}[subsection]{Remarks}

\numberwithin{equation}{subsection}

% Macros
%
\def\lim{\mathop{\mathrm{lim}}\nolimits}
\def\colim{\mathop{\mathrm{colim}}\nolimits}
\def\Spec{\mathop{\mathrm{Spec}}}
\def\Hom{\mathop{\mathrm{Hom}}\nolimits}
\def\Ext{\mathop{\mathrm{Ext}}\nolimits}
\def\SheafHom{\mathop{\mathcal{H}\!\mathit{om}}\nolimits}
\def\SheafExt{\mathop{\mathcal{E}\!\mathit{xt}}\nolimits}
\def\Sch{\mathit{Sch}}
\def\Mor{\mathop{\mathrm{Mor}}\nolimits}
\def\Ob{\mathop{\mathrm{Ob}}\nolimits}
\def\Sh{\mathop{\mathit{Sh}}\nolimits}
\def\NL{\mathop{N\!L}\nolimits}
\def\CH{\mathop{\mathrm{CH}}\nolimits}
\def\proetale{{pro\text{-}\acute{e}tale}}
\def\etale{{\acute{e}tale}}
\def\QCoh{\mathit{QCoh}}
\def\Ker{\mathop{\mathrm{Ker}}}
\def\Im{\mathop{\mathrm{Im}}}
\def\Coker{\mathop{\mathrm{Coker}}}
\def\Coim{\mathop{\mathrm{Coim}}}

% Boxtimes
%
\DeclareMathSymbol{\boxtimes}{\mathbin}{AMSa}{"02}

%
% Macros for moduli stacks/spaces
%
\def\QCohstack{\mathcal{QC}\!\mathit{oh}}
\def\Cohstack{\mathcal{C}\!\mathit{oh}}
\def\Spacesstack{\mathcal{S}\!\mathit{paces}}
\def\Quotfunctor{\mathrm{Quot}}
\def\Hilbfunctor{\mathrm{Hilb}}
\def\Curvesstack{\mathcal{C}\!\mathit{urves}}
\def\Polarizedstack{\mathcal{P}\!\mathit{olarized}}
\def\Complexesstack{\mathcal{C}\!\mathit{omplexes}}
% \Pic is the operator that assigns to X its picard group, usage \Pic(X)
% \Picardstack_{X/B} denotes the Picard stack of X over B
% \Picardfunctor_{X/B} denotes the Picard functor of X over B
\def\Pic{\mathop{\mathrm{Pic}}\nolimits}
\def\Picardstack{\mathcal{P}\!\mathit{ic}}
\def\Picardfunctor{\mathrm{Pic}}
\def\Deformationcategory{\mathcal{D}\!\mathit{ef}}
\usepackage{cleveref}
\crefname{originalreference}{original reference}{original references}
\Crefname{originalreference}{Original reference}{Original references}
\makeatletter
\newlabel{lemma-finite-directed-colimit-surjective-maps}{{086M}{source}{}{sourcepointer-1}{}}
\newlabel{lemma-finite-directed-colimit-surjective-maps@cref}{{[originalreference][0][]086M}{[1][0][]source}{}{}{}}
\makeatother
\crefname{originalreference}{original reference}{original references}
\Crefname{originalreference}{Original reference}{Original references}
\makeatletter
\newlabel{lemma-quasi-coherent-colimit-finite-type}{{01PG}{source}{}{sourcepointer-2}{}}
\newlabel{lemma-quasi-coherent-colimit-finite-type@cref}{{[originalreference][0][]01PG}{[1][0][]source}{}{}{}}
\makeatother
\title{A finite quasi-coherent algebra is a directed colimit of finite finitely presented algebras with surjective transition maps}
\author{The Stacks Project Authors}
\date{stacks-project:properties.tex:local-snapshot-20261012; GNU FDL 1.2 or later}
\begin{document}
\maketitle
\noindent\textit{Editorial note.} All delivered mathematical text is the author's own, copied from the pinned Stacks Project chapter file properties.tex; no statement or proof step was rewritten, reordered or removed.\par\medskip
\noindent\textit{Editorial note.} The counted proof is the author's complete argument: the presentation of the algebra as a quotient of the symmetric algebra on a finitely presented module, the approximation of the kernel by finite type submodules giving finitely presented approximants with surjective transition maps, and the monic-relation argument making the approximants finite for large index.\par\medskip
\noindent\textit{Editorial note.} The label lemma-finite-algebra-directed-colimit-finite-finitely-presented names the counted statement, and the proof environment immediately following it in the source is delivered verbatim as its proof.\par\medskip
\noindent\textit{Editorial note.} The counted statement is carried by the author's own lemma environment in the chapter source.\par\medskip
\noindent\textit{Editorial note.} Pre-open gate: python operations/agent-authoring/tools/precard\_gate.py --source properties.tex --label lemma-finite-algebra-directed-colimit-finite-finitely-presented --start 3597 --end 3638 --local-id 28037 returned RESULT=OK (L28037 free, occupancy FREE, public interval clear). Candidate taken from sources/stacks-properties/candidates-halfA.json.\par\medskip
\noindent\textit{Editorial note.} The statement and proof are copied from the pinned chapter file; the cited inputs are kept as external pointers with their canonical Stacks tags.\par\medskip
\noindent\textit{Allowed standard prerequisites (editorial).}\begin{itemize}
\item Presentation of a finite quasi-coherent algebra by a finitely presented module. Same chapter, the lemma producing a finitely presented module with a surjection onto the given quasi-coherent module. Supplies the finitely presented module F with a surjection onto the algebra A.
\item Description of a quasi-coherent module as a filtered colimit of finite type submodules. Same chapter, the lemma writing a quasi-coherent module as a filtered colimit of its finite type submodules. Supplies the presentation of the kernel of the symmetric algebra surjection as a filtered colimit of finite type submodules.
\item Integrality of a finite algebra over an affine open. Algebra, the characterization of a finite algebra over a ring through monic polynomial relations for generators. Supplies the monic polynomials satisfied by the chosen generators over each affine open.
\end{itemize}\medskip
\ifdefined\hypertarget\hypertarget{sourcepointer-1}{}\fi
\noindent\textit{Source reference.} Original reference 086M: sources/stacks-properties/pinned-original-source.tex, source lines 3442--3442. Prerequisite pointer invoked at proof line(s) 3608; it names a standard fact of the same chapter that the counted argument uses, and every step of the counted proof built on it is retained verbatim in the delivered body.\par\medskip
\ifdefined\hypertarget\hypertarget{sourcepointer-2}{}\fi
\noindent\textit{Source reference.} Original reference 01PG: sources/stacks-properties/pinned-original-source.tex, source lines 3115--3115. Prerequisite pointer invoked at proof line(s) 3617; it names a standard fact of the same chapter that the counted argument uses, and every step of the counted proof built on it is retained verbatim in the delivered body.\par\medskip
\begin{lemma}
\label{lemma-finite-algebra-directed-colimit-finite-finitely-presented}
Let $X$ be a scheme. Assume $X$ is quasi-compact and quasi-separated.
Let $\mathcal{A}$ be a finite quasi-coherent $\mathcal{O}_X$-algebra.
Then $\mathcal{A} = \colim \mathcal{A}_i$ is a directed colimit of finite
and finitely presented quasi-coherent $\mathcal{O}_X$-algebras
such that all transition maps $\mathcal{A}_{i'} \to \mathcal{A}_i$
are surjective.
\end{lemma}

\begin{proof}
By Lemma \ref{lemma-finite-directed-colimit-surjective-maps}
there exists a finitely presented $\mathcal{O}_X$-module
$\mathcal{F}$ and a surjection $\mathcal{F} \to \mathcal{A}$.
Using the algebra structure we obtain a surjection
$$
\text{Sym}^*_{\mathcal{O}_X}(\mathcal{F}) \longrightarrow \mathcal{A}
$$
Denote $\mathcal{J}$ the kernel. Write $\mathcal{J} = \colim \mathcal{E}_i$
as a filtered colimit of finite type $\mathcal{O}_X$-submodules
$\mathcal{E}_i$ (Lemma \ref{lemma-quasi-coherent-colimit-finite-type}). Set
$$
\mathcal{A}_i = \text{Sym}^*_{\mathcal{O}_X}(\mathcal{F})/(\mathcal{E}_i)
$$
where $(\mathcal{E}_i)$ indicates the ideal sheaf generated by
the image of $\mathcal{E}_i \to \text{Sym}^*_{\mathcal{O}_X}(\mathcal{F})$.
Then each $\mathcal{A}_i$ is a finitely presented $\mathcal{O}_X$-algebra,
the transition maps are surjections,
and $\mathcal{A} = \colim \mathcal{A}_i$. To finish the proof we still
have to show that $\mathcal{A}_i$ is a finite $\mathcal{O}_X$-algebra
for $i$ sufficiently large. To do this we choose an affine open
covering $X = U_1 \cup \ldots \cup U_m$. Take generators
$f_{j, 1}, \ldots, f_{j, N_j} \in \Gamma(U_i, \mathcal{F})$.
As $\mathcal{A}(U_j)$ is a finite $\mathcal{O}_X(U_j)$-algebra we
see that for each $k$ there exists a monic polynomial
$P_{j, k} \in \mathcal{O}(U_j)[T]$ such that $P_{j, k}(f_{j, k})$
is zero in $\mathcal{A}(U_j)$. Since
$\mathcal{A} = \colim \mathcal{A}_i$ by construction, we
have $P_{j, k}(f_{j, k}) = 0$ in $\mathcal{A}_i(U_j)$
for all sufficiently large $i$. For such $i$ the algebras
$\mathcal{A}_i$ are finite.
\end{proof}

\end{document}
